% Source: https://artofproblemsolving.com/community/c3h2472984_is_0000001__0
% Stable identity: aops-topic-2472984-q01
\textbf{(a)} A decimal expansion assigns one digit to each position $1,2,3,\dots$ after the point. In $0.\overline{0}1$ every position $k$ holds a $0$ (the repeating block fills all positions), so there is no position left for the $1$: no finite position comes ``after infinitely many zeros''. The expansion with a $0$ in every position is just $0.000\ldots=0$. What the student probably has in mind is the sequence $0.1,\ 0.01,\ 0.001,\dots$, that is $10^{-n}$, whose terms get as close to $0$ as we like; that is what (b) makes precise.

\textbf{(b)} Suppose $x>0$. Since $10^n\ge n+1$ for every positive integer $n$ (by induction, or Bernoulli's inequality), we can choose $n$ with $10^n>\dfrac1x$, i.e. $x>10^{-n}$. This contradicts $x<10^{-n}$. Hence $x\le0$, and with $x\ge0$ we get $\boxed{x=0}$. $\blacksquare$
